\section{Irreducibility}\label{s:irreducible}

In this section, we study the irreducibility condition.
It will be convenient to work in the \defn{augmented simplex category}, which contains an additional object \([-1] = \emptyset\) and a unique morphism from it to every \([n]\).
The natural equivalence between $\cP$ and $\chains$ of \cref{st:chains_functor_equivalence} extends to this category; we denote the extended functors by $\widehat{\cP}$ and $\widehat\chains$, respectively.
For every $n$, the only additional basis element of $\widehat{\cP}(\simplex^n)$ is the full subset $\set{0,\dots,n}$ in degree $-1$.
We freely identify the underlying graded vector spaces of $\widehat{\cP}(\simplex^n)$ and $\cP(\simplex^n)$ in non-negative degrees.

The main results of this section are an explicit description, in terms of a Hopf algebra structure carried by $\widehat{\cP}(\simplex^n)$, of the largest subcomplex of irreducible chains, and a proof that irreducibility together with semi-simpliciality implies simpliciality.

\subsection{A Koszul model}

Recall that an element $a$ of a commutative $\F$-algebra $A$ with $a^2 = 0$ makes $(A,\, a \cdot -)$ into a chain complex, referred to as the \defn{Koszul complex} of $a$.
Consider the exterior algebra over $\F$ generated by $x_0, \dots, x_n$, one generator for each vertex of $\simplex^n$, graded so that the monomial $x_U$ associated to $U \subseteq \set{0,\dots,n}$ has degree $n - \bars{U}$, with $x_\emptyset = 1$.
The assignment $U \mapsto x_U$ defines an isomorphism of complexes between $\widehat{\cP}(\simplex^n)$ and the Koszul complex of
\[
\omega = x_0 + \dots + x_n.
\]
We use it to regard $\widehat{\cP}(\simplex^n)$ as an algebra.
%since $\omega\, x_U = \sum_{\bar u \notin U} x_{\bar u . U}$, the terms with $\bar u \in U$ vanishing because $x_{\bar u}^2 = 0$.
In what follows we use $U$ and $x_U$ interchangeably.
This Koszul complex factors as the tensor product of the complexes $\big(\F[x_u]/(x_u^2),\, x_u \cdot -\big)$, each of them acyclic.

\begin{remark*}\label{r:group_ring}
	Writing $t_u = 1 + x_u$ identifies the algebra $\widehat{\cP}(\simplex^n)$ with the group ring $\F[C_2^{\, n+1}]$ of an elementary abelian $2$-group, under which the Koszul element $\omega = \sum_u (1 + t_u)$ is the sum of the augmentation elements of its factors.
	The differentials of the complex $W$ of \cref{s:statement} are of the same nature, being multiplication by the augmentation element $1 + T$ of $\F[C_2]$.
\end{remark*}

\subsection{The partition coproduct}

The exterior algebra is a Hopf algebra, with coproduct the algebra map determined by declaring the generators primitive: $\Delta(x_u) = x_u \ot 1 + 1 \ot x_u$.
Explicitly,
\[
\Delta(x_U) = \sum x_V \ot x_W,
\]
the sum taken over all $V$ and $W$ satisfying $V \union W = U$ and $V \cap W = \emptyset$.
We refer to $\zeta_U = \Delta(x_U) \in \widehat{\cP}(\simplex^n)^{\ot 2}$ as the \defn{partition chain} of $U$.
Since $\Delta$ is cocommutative, partition chains are fixed by $T$.
If $U \subsetneq \set{0,\dots,n}$, then $\zeta_U$ is concentrated in non-negative bidegrees and agrees with its image in the ordinary complex $\cP(\simplex^n)^{\ot 2}$.
With respect to the grading above, $\Delta$ raises degree by $n$, so it defines a map $\Sigma^n \widehat{\cP}(\simplex^n) \to \widehat{\cP}(\simplex^n)^{\ot 2}$ from the $n$-fold suspension.
The next lemma states that this map is a chain map.

\begin{lemma}\label{c:boundary_partition_chain}
	For every $U \subseteq \set{0,\dots,n}$ we have
	\begin{equation}\label{eq:boundary_partition_chain}
		\bd\zeta_U = \sum_{\bar u \notin U} \zeta_{\bar u.U}.
	\end{equation}
\end{lemma}

\begin{proof}
	The differential of $\widehat{\cP}(\simplex^n)^{\ot 2}$ is multiplication by $\Delta(\omega) = \omega \ot 1 + 1 \ot \omega$.
	Since $\Delta$ is an algebra map,
	\[
	\bd \Delta(x_U) = \Delta(\omega)\, \Delta(x_U) = \Delta(\omega\, x_U) = \sum_{\bar u \notin U} \Delta(x_{\bar u.U}). \qedhere
	\]
\end{proof}

The proof exhibits $\bd$ as a coderivation of $\Delta$, reflecting the primitivity of $\omega$.
It is, however, not a derivation of the product: $\bd(x_V x_W) = \omega\, x_V x_W$, whereas $\bd(x_V)\, x_W + x_V\, \bd(x_W) = 2\, \omega\, x_V x_W = 0$.
Thus, $\widehat{\cP}(\simplex^n)$ is a differential graded coalgebra but not a differential graded Hopf algebra.

\begin{remark*}
	The partition coproduct does \emph{not} correspond to the group-ring coproduct $\F[C_2^{\, n+1}]$ as it makes each $x_u$ primitive rather than each $t_u$ grouplike.
%	The two are non-isomorphic as Hopf algebras since the group ring has $2^{n+1}$ grouplike elements, whereas the partition coproduct has only one.
%	The primitively generated structure is the Cartier dual of the group ring, and also the mod~$2$ homology of the $(n+1)$-torus.
\end{remark*}


\subsection{Irreducible chains}

A basis element $x_V \ot x_W \in \widehat{\cP}(\simplex^n)^{\ot 2}$ is said to be \defn{irreducible} if $V \cap W = \emptyset$ (equivalently, if the product $x_Vx_W$ is non-zero) and \defn{reducible} otherwise.
We denote by $\pired$ and $\piirred$ the projections onto the graded subspaces generated by reducible and irreducible basis elements, respectively.
Under the identification above, a cup-$i$ construction $\triangle$ is irreducible if and only if its defining elements $\triangle_i[n]$ are irreducible in this sense for all $i,n \in \N$.

\begin{remark}
	Irreducibility admits a geometric interpretation.
	Under the equivalence of \cref{st:chains_functor_equivalence}, a basis element $x_V \ot x_W$ corresponds to the pair of faces $\big(d_V[n],\, d_W[n]\big)$, and it is irreducible precisely when
	\[
	\dim d_V[n] + \dim d_W[n] - \dim \big(d_V[n] \cap d_W[n]\big) = n,
	\]
	that is, when the two faces intersect transversally in $\simplex^n$, in which case the product $x_V x_W$ corresponds to their intersection $d_{V \union W}[n]$.
	The multiplication of $\widehat{\cP}(\simplex^n)$ is thus the mod~$2$ intersection product of transverse faces, and the irreducibility axiom states that a \mbox{cup-$i$} construction is supported on transverse pairs.
\end{remark}

The reducible basis elements span the ideal of $\widehat{\cP}(\simplex^n)^{\ot 2}$ generated by the elements $x_k \ot x_k$.
Since the differential of $\widehat{\cP}(\simplex^n)^{\ot 2}$ is multiplication by $\Delta(\omega)$, every ideal is closed under it; in particular, the reducible subspace is a subcomplex.
We remark that this ideal is properly contained in the kernel of the multiplication map $\widehat{\cP}(\simplex^n)^{\ot 2} \to \widehat{\cP}(\simplex^n)$, which also contains irreducible chains such as $x_i \ot x_j + x_j \ot x_i$.
The irreducible subspace is not a subcomplex, and we now describe the largest subcomplex it contains.

\begin{theorem}\label{st:augmented_partition_subcomplex}
	Partition chains span the largest subcomplex contained in the irreducible subspace of $\widehat{\cP}(\simplex^n)^{\ot 2}$.
\end{theorem}

\begin{proof}
	Let $I$ be the irreducible subspace of $\widehat{\cP}(\simplex^n)^{\ot 2}$ and let $A$ be the span of the partition chains.
	The largest subcomplex contained in \(I\) is \(I \cap \bd^{-1}(I)\).
	The summands of every partition chain are irreducible, so $A \subseteq I$, and $A$ is a subcomplex by \cref{c:boundary_partition_chain}; hence \(A \subseteq I \cap \bd^{-1}(I)\).
	For the other inclusion,
%	Since any subcomplex contained in $I$ is contained in $I \cap \bd^{-1}(I)$, it suffices to prove that
%	\begin{equation}\label{eq:augmented_irreducible_kernel}
%		I \cap \bd^{-1}(I) \subseteq A.
%	\end{equation}
%	For the other,
	let $\zeta = \sum_\Lambda V \ot W$ be in $I \cap \bd^{-1}(I)$, where $\Lambda$ is a set of distinct irreducible basis elements.
	The reducible part of the boundary of $\zeta$ is
	\begin{equation}\label{eq:reducible_augmented_boundary}
		\pired(\bd\zeta) =
		\sum_\Lambda\left(
		\sum_{w \in W} w.V \ot W +
		\sum_{v \in V} V \ot v.W
		\right).
	\end{equation}
	Distinct terms within each inner family are distinct, while a term of the first family agrees with one of the second precisely when the corresponding partitions differ by a \defn{basic move}: the transfer of a single element between the factors.
	Hence \cref{eq:reducible_augmented_boundary} vanishes if and only if $\Lambda$ is closed under basic moves.
	Since basic moves preserve unions and connect any two partitions of the same set, this occurs exactly when $\Lambda$ is a union of complete families of partitions, that is, when $\zeta \in A$.
\end{proof}

We will apply this result to chains in the ordinary complex.
If $\zeta \in \cP(\simplex^n)^{\ot 2}$ is a homogeneous irreducible chain of degree at least $n$, then the augmented and ordinary boundaries of $\zeta$ have the same reducible part: their difference consists of terms having the full subset as one factor, the degree assumption forces the other factor to be $\emptyset$, and such terms are irreducible.
Consequently, if $\pired(\bd\zeta) = 0$ then $\zeta$ is a sum of partition chains, necessarily indexed by sets with at most $n$ elements.

\subsection{Irreducibility implies simpliciality}

The following theorem is not used in this paper but it is included to clarify the relationship between two axioms: irreducibility and naturality.

\begin{theorem}\label{t:irreducible_implies_simplicial}
	Every irreducible semi-simplicial \mbox{cup-$i$} construction is simplicial.
\end{theorem}

\begin{proof}
	Let $\triangle$ be an irreducible semi-simplicial \mbox{cup-$i$} construction.
	By \cref{l:simplicial_from_semisimplicial}, it suffices to show that
	\[
		\cP(\sigma_j)^{\ot 2}\triangle_i[n] = 0
	\]
	for every codegeneracy $\sigma_j \colon [n] \to [n-1]$.
	We derive this from the reducible part of the cup-$i$ identity.
	Fix $i,n \in \N$ with $n > 0$.
	We use the coefficient notation $\coeff{\zeta}{V \ot W}$ introduced in \cref{ss:axioms_revisited} and remind the reader that each $\triangle_i[n]$ has degree $n+i\geq n$.

%	, so the preceding comparison allows us to compute the reducible part of its boundary in the augmented complex while retaining the ordinary cup-$i$ identities.

	We first claim that
	\begin{equation}\label{eq:reducible_part}
		\pired\big(\bd\triangle_i[n]\big) = \triangle_i\bd\,[n].
	\end{equation}
	Indeed, $\triangle_i\bd\,[n] = \sum_{k=0}^{n} \cP(\delta_k)^{\ot 2}\triangle_i[n-1]$ and, since every subset in the image of $\cP(\delta_k)$ contains $k$, the basis elements in the image of $\cP(\delta_k)^{\ot 2}$ are reducible.
	On the other hand, $(1+T)\triangle_{i-1}[n]$ is a sum of irreducible basis elements.
	Applying $\pired$ to the identity $\bd\triangle_i[n] + \triangle_i\bd\,[n] = (1+T)\triangle_{i-1}[n]$ proves \eqref{eq:reducible_part}.

	Let us compare coefficients on both sides of \eqref{eq:reducible_part}.
	By \cref{eq:reducible_augmented_boundary}, every basis element appearing on the left-hand side of \cref{eq:reducible_part} is of the form $X \ot Y$ with $X \cap Y = \set{k}$ for exactly one $k$, and the same holds on the right-hand side, since $\cP(\delta_k)$ adjoins $k$ to both factors of an irreducible basis element.
	Fix such an $X \ot Y$.
	Its coefficient on the left-hand side is
	\[
	\coeff{\triangle_i[n]}{X \setminus k \ot Y} + \coeff{\triangle_i[n]}{X \ot Y \setminus k}.
	\]
	For the right-hand side, notice that $\cP(\delta_m)^{\ot 2}(A \ot B) = X \ot Y$ forces $m = k$, and that $\cP(\delta_k)$ defines a bijection from the subsets of $\set{0,\dots,n-1}$ to the subsets of $\set{0,\dots,n}$ containing $k$, with inverse $S \mapsto \sigma_k(S \setminus k)$, where $\sigma_k$ is the elementwise relabeling of \cref{eq:codegeneracies}.
	Therefore, \cref{eq:reducible_part} is equivalent to the identity
	\begin{equation}\label{eq:coefficient_comparison}
		\coeff{\triangle_i[n]}{X \setminus k \ot Y} + \coeff{\triangle_i[n]}{X \ot Y \setminus k} =
		\coeff{\triangle_i[n-1]}{\sigma_k(X \setminus k) \ot \sigma_k(Y \setminus k)}
	\end{equation}
	holding for every basis element $X \ot Y$ with $X \cap Y = \set{k}$.

	Let $\tau$ denote the exchange of $j$ and $j+1$, acting on subsets of $\set{0,\dots,n}$.
	Inspecting \cref{eq:codegeneracies}, we see that $\cP(\sigma_j)(\tau U) = \cP(\sigma_j)(U)$ for every subset $U$, and that an irreducible basis element $V \ot W$ satisfies $\cP(\sigma_j)^{\ot 2}(V \ot W) \neq 0$ if and only if one of $j,j+1$ belongs to $V$ and the other belongs to $W$.
	Such basis elements come in pairs $\set[\big]{V \ot W,\ \tau(V) \ot \tau(W)}$ of distinct elements sharing the same non-zero image, and elements of different pairs have different images.
	Consequently, since $\triangle_i[n]$ is irreducible, $\cP(\sigma_j)^{\ot 2}\triangle_i[n] = 0$ if and only if
	\begin{equation}\label{eq:swap}
		\coeff{\triangle_i[n]}{V \ot W} = \coeff{\triangle_i[n]}{\tau(V) \ot \tau(W)}
	\end{equation}
	for every such element.

	To establish \cref{eq:swap}, assume that $j \in V$ and $j+1 \in W$, the other case being symmetric.
	Since $W$ contains $j+1$ but not $j$, the set $W \union \set{j}$ is fixed by $\tau$ and equals $\tau(W) \union \set{j{+}1}$, so the basis elements $V \ot \big(W \union \set{j}\big)$ and $\tau(V) \ot \big(W \union \set{j}\big)$ have factors intersecting in $\set{j}$ and $\set{j+1}$, respectively.
	Applying \cref{eq:coefficient_comparison} to each of them gives
	\begin{align*}
		\coeff{\triangle_i[n]}{V \setminus j \ot W \union \set{j}} &+ \coeff{\triangle_i[n]}{V \ot W} \\ &=
		\coeff{\triangle_i[n-1]}{\sigma_j(V \setminus j) \ot \sigma_j(W)}, \\
		\coeff{\triangle_i[n]}{V \setminus j \ot W \union \set{j}} &+ \coeff{\triangle_i[n]}{\tau(V) \ot \tau(W)} \\ &=
		\coeff{\triangle_i[n-1]}{\sigma_{j+1}(V \setminus j) \ot \sigma_{j+1}\big(\tau(W)\big)}.
	\end{align*}
	The maps $\sigma_j$ and $\sigma_{j+1}$ agree on subsets containing neither $j$ nor $j+1$, so the first factors on the right agree, and so do the second, both being $\sigma_j\big(W \setminus \set{j{+}1}\big) \union \set{j}$.
	Adding the two identities therefore yields \cref{eq:swap}.
\end{proof}
